% GENERATED FILE. Edit canonical lessons, then run npm run generate:books.
% canonical-source-hash: fnv1a64:a4a4b86c9185ff2c
% canonical-lessons: TE-C287-chati, TE-C287-pramadam, TE-C287-ambulensu, TE-C287-bima, TE-C287-atyavasaram

\chapter{Chest, Accident, Ambulance, Insurance, Emergency}
\label{ch:te-chest-accident-ambulance-insurance-emergency}
\hlchaptermodality{287}

\begin{chapteropening}
\textbf{By the end of this chapter:} \emph{I can say the chest, an accident, an ambulance, insurance and an emergency.}

Complete the last lesson of chapter 287: I can say the chest, an accident, an ambulance, insurance and an emergency.
\end{chapteropening}

\section[\texorpdfstring{\te{ఛాతీ} (chātī)}{ఛాతీ (chātī)}]{\te{ఛాతీ} (chātī) — the chest}

\label{lesson:TE-C287-chati}



\begin{quote}
Before the new one: say the Telugu for to drain, then the Telugu for to lean.
\end{quote}

\subsection*{You'll want to know: \te{ఛాతీ}}
\textbf{\te{ఛాతీ}} — \emph{chātī} — \textquotedblleft{}the chest\textquotedblright{}.

\begin{grammarlens}[title={What you've built}]
One more word for this chapter.
\end{grammarlens}

\subsection*{Your turn}
\begin{itemize}
\raggedright
  \item \emph{Say it:} \emph{chātī}
  \item \emph{Say it:} \emph{chātī}, once more
  \item \emph{Say it:} say \emph{chātī}
  \item \emph{Recall:} say the Telugu for to drain, then the Telugu for to lean, then say \emph{chātī} again
\end{itemize}

\begin{tcolorbox}[breakable,colback=teal!4,colframe=teal!35!black,title={Before you move on}]
What does \te{ఛాతీ} mean? (\textquotedblleft{}The chest\textquotedblright{}.) Say it once more.
\end{tcolorbox}
\section[\texorpdfstring{\te{ప్రమాదం} (pramādaṁ)}{ప్రమాదం (pramādaṁ)}]{\te{ప్రమాదం} (pramādaṁ) — an accident; danger}

\label{lesson:TE-C287-pramadam}



\begin{quote}
Before the new one: say the Telugu for to lean, then the Telugu for the chest.
\end{quote}

\subsection*{You'll want to know: \te{ప్రమాదం}}
\textbf{\te{ప్రమాదం}} — \emph{pramādaṁ} — \textquotedblleft{}an accident; danger\textquotedblright{}.

\begin{grammarlens}[title={What you've built}]
One more word for this chapter.
\end{grammarlens}

\subsection*{Your turn}
\begin{itemize}
\raggedright
  \item \emph{Say it:} \emph{pramādaṁ}
  \item \emph{Say it:} \emph{pramādaṁ}, once more
  \item \emph{Say it:} say \emph{pramādaṁ}
  \item \emph{Recall:} say the Telugu for to lean, then the Telugu for the chest, then say \emph{pramādaṁ} again
\end{itemize}

\begin{tcolorbox}[breakable,colback=teal!4,colframe=teal!35!black,title={Before you move on}]
What does \te{ప్రమాదం} mean? (\textquotedblleft{}An accident; danger\textquotedblright{}.) Say it once more.
\end{tcolorbox}
\section[\texorpdfstring{\te{అంబులెన్సు} (ambulensu)}{అంబులెన్సు (ambulensu)}]{\te{అంబులెన్సు} (ambulensu) — an ambulance}

\label{lesson:TE-C287-ambulensu}



\begin{quote}
Before the new one: say the Telugu for the chest, then the Telugu for an accident.
\end{quote}

\subsection*{You'll want to know: \te{అంబులెన్సు}}
\textbf{\te{అంబులెన్సు}} — \emph{ambulensu} — \textquotedblleft{}an ambulance\textquotedblright{}.

\begin{grammarlens}[title={What you've built}]
One more word for this chapter.
\end{grammarlens}

\subsection*{Your turn}
\begin{itemize}
\raggedright
  \item \emph{Say it:} \emph{ambulensu}
  \item \emph{Say it:} \emph{ambulensu}, once more
  \item \emph{Say it:} say \emph{ambulensu}
  \item \emph{Recall:} say the Telugu for the chest, then the Telugu for an accident, then say \emph{ambulensu} again
\end{itemize}

\begin{tcolorbox}[breakable,colback=teal!4,colframe=teal!35!black,title={Before you move on}]
What does \te{అంబులెన్సు} mean? (\textquotedblleft{}An ambulance\textquotedblright{}.) Say it once more.
\end{tcolorbox}
\section[\texorpdfstring{\te{బీమా} (bīmā)}{బీమా (bīmā)}]{\te{బీమా} (bīmā) — insurance}

\label{lesson:TE-C287-bima}



\begin{quote}
Before the new one: say the Telugu for an accident, then the Telugu for an ambulance.
\end{quote}

\subsection*{You'll want to know: \te{బీమా}}
\textbf{\te{బీమా}} — \emph{bīmā} — \textquotedblleft{}insurance\textquotedblright{}.

\begin{grammarlens}[title={What you've built}]
One more word for this chapter.
\end{grammarlens}

\subsection*{Your turn}
\begin{itemize}
\raggedright
  \item \emph{Say it:} \emph{bīmā}
  \item \emph{Say it:} \emph{bīmā}, once more
  \item \emph{Say it:} say \emph{bīmā}
  \item \emph{Recall:} say the Telugu for an accident, then the Telugu for an ambulance, then say \emph{bīmā} again
\end{itemize}

\begin{tcolorbox}[breakable,colback=teal!4,colframe=teal!35!black,title={Before you move on}]
What does \te{బీమా} mean? (\textquotedblleft{}Insurance\textquotedblright{}.) Say it once more.
\end{tcolorbox}
\section[\texorpdfstring{\te{అత్యవసరం} (atyavasaraṁ)}{అత్యవసరం (atyavasaraṁ)}]{\te{అత్యవసరం} (atyavasaraṁ) — an emergency; urgent}

\label{lesson:TE-C287-atyavasaram}



\begin{quote}
Before the new one: say the Telugu for an ambulance, then the Telugu for insurance.
\end{quote}

\subsection*{You'll want to know: \te{అత్యవసరం}}
\textbf{\te{అత్యవసరం}} — \emph{atyavasaraṁ} — \textquotedblleft{}an emergency; urgent\textquotedblright{}.

\begin{grammarlens}[title={What you've built}]
One more word for this chapter.
\end{grammarlens}

\subsection*{Your turn}
\begin{itemize}
\raggedright
  \item \emph{Say it:} \emph{atyavasaraṁ}
  \item \emph{Say it:} \emph{atyavasaraṁ}, once more
  \item \emph{Say it:} say \emph{atyavasaraṁ}
  \item \emph{Recall:} say the Telugu for an ambulance, then the Telugu for insurance, then say \emph{atyavasaraṁ} again
\end{itemize}

\begin{tcolorbox}[breakable,colback=teal!4,colframe=teal!35!black,title={Before you move on}]
What does \te{అత్యవసరం} mean? (\textquotedblleft{}An emergency; urgent\textquotedblright{}.) Say it once more.
\end{tcolorbox}
